\heading{数据结构与算法A-14秋-期末}
\noteinfo{题目来源为真题扫描件转录。客观题答案和部分参考解答已整理在解答中；原扫描中考场纪律、装订线等非题目信息已尽量删去。本次整理提供 C 语言版与 Python 语言版；除程序代码语言外，两版题目内容保持一致。}

\begin{question}
一、填空（27 分）\\
1. （2 分）当各边上的权值 \underline{\hspace{2em}} 时，BFS 算法可用来解决单源最短路径问题。\\
A．均相等        B．均互不相等      C．不一定相等\\
2. （4 分）有向图G 如下图所示：\
\pyfig[0.32\linewidth]{figures/数据结构与算法A-14秋-期末-fig01.png}{}
\
（1）写出所有可能的拓扑序列： \underline{\hspace{8em}}。\\
（2）添加一条弧 \underline{\hspace{8em}} 之后，则仅有唯一的拓扑序列。\\
3. （2 分）设有如下所示无向图G，用Prim 算法构造最小生成树所走过的边\
\pyfig[0.36\linewidth]{figures/数据结构与算法A-14秋-期末-fig02.png}{}
\
的集合E=\{(1，3)，(1，2)，(3，5)，(5，6)，(6，4)\} （用边（2，3）代替\\
（1，2）也可以）。\\
4. （2 分）如果只想得到1000 个元素组成的序列中第5 个最小元素之前的部\\
分排序的序列，用 \underline{\hspace{2em}} 方法最快。\\
A．起泡排序   B．快速排列   C．堆排序  D．简单选择排序\\
5. （2 分）若要求尽可能快地对序列进行稳定的排序，则应选 \underline{\hspace{2em}}\\
A．快速排序       B．归并排序       C．冒泡排序       D.归并排序\\
V1\\
V3\\
V4\\
V2\\
6. （2 分）在文件局部有序或文件长度较小的情况下，最佳的排序方法是 \underline{\hspace{2em}}。\\
A.直接插入排序    B. 冒泡排序       C. 直接选择排序    D.归并排序\\
7. （2 分）设输入的关键字满足K1>K2> … >Kn ，缓冲区大小为 m ，用置换\\
选择排序方法可产生 \underline{\hspace{4em}} 个初始归并段。\\
8. （2 分）设有13 个初始归并段，其长度分别为28，16，37，42，5，9，13，\\
14，20，17，30，12，18。试计算最佳4 路归并树的带权路径长度WPL= \underline{\hspace{4em}}。\\
9. （3 分）设有序顺序表中的元素依次为017, 094, 154, 170, 275, 503, 509, 512,\\
553, 612, 677, 765, 897, 908，则对其进行折半查找时，等概率下搜索成功的\\
平均查找长度和不成功的平均查找长度分别为 \underline{\hspace{4em}} 和 \underline{\hspace{4em}}。\\
10. （2 分）设散列表长度为14（地址下标为0…13），哈希函数为H(key)=key\%11，\\
表中已有数据的关键字为15,38,61,84 共四个，现要将关键字为49 的结点将\\
加入表中，用二次探查法解决冲突，则放入的位置是 \underline{\hspace{2em}}。\\
A．8           B．3           C．5            D．9\\
11. （2 分）在一棵含有1023 个关键字的4 阶B 树中进行查找（不读取数据主\\
文件），至多读盘 \underline{\hspace{2em}} 次。\\
A. 7           B. 8            C.  10          D. 9\\
12. （2 分）在一棵阶为 k 的红黑树中，内部结点最少有2k-1 个；从根到叶的\\
简单路径长度的范围为 [k,2k]。\\
二、 辨析与简答（30 分）\\
1.\\
（8分）将关键字序列（7、8、30、11、18、9、14）散列存储到散列表中。\\
散列表是一个下标从0开始的一维数组，散列函数为：H(key) = (key * 3) MOD\\
7，处理冲突采用线性探测法，要求装填（载）因子为0.7。\\
(1) 请画出所构造的散列表。\\
(2) 分别计算等概率情况下查找成功和查找不成功的平均查找长度。\\
2.\\
（6分）设有如下一棵3阶B树，请画出在其中插入关键码19后的B树，并给出\
\pyfig[0.70\linewidth]{figures/数据结构与算法A-14秋-期末-fig03.png}{}
\
插入过程中的访外（读写）次数；在此基础上再删除关键码48，请画出删除\\
后的B树及删除过程中的访外（读写）次数。\\
3. （6分）将下列字符序列 E A S Y Q U E S T I O N 依次插入到初始为空的红黑\\
树（RB-tree）中，请画出最终得到的红黑树。（用圆圈表示红色结点，方框表\\
示黑色结点，外部空叶结点可省略不画）。\\
4. （4分）利用广义表的Head 和 Tail运算，把原子d分别从下列广义表中分离出来，\\
L1＝(((((a),b),d),e))；L2＝（a,(b,((d)),e））。\\
5. （6分）将关键码1，2，3，...，（2k-1）依次插入到一棵初始为空的AVL树中，\\
试证明结果是一棵高度为k的完全满二叉树。\\
三、算法填空 （18 分）\\
阅读和完成下面的代码（每空可不限一条语句）。\\
1.\\
（10 分）以单向链表为存储结构，实现选择排序算法（得到非递减序列）。\\
class node \{\\
int key;\\
node *next;\\
\};\\
node* sort\_linked\_list(node* head) \{\\
i = head;\\
while (i!=None) \{\\
填空1 \underline{\hspace{12em}};\\
while (j!=None) \{\\
if (填空2 \underline{\hspace{12em}})   填空3 \underline{\hspace{12em}};\\
填空4 \underline{\hspace{12em}};\\
\}\\
填空5 \underline{\hspace{12em}};\\
\}\\
return head;\\
\}\\
2.\\
（8 分）下面是败者树在内部结点从右分支向上比赛的成员函数实现，请将空\\
缺部分补充完整。\\
template<class T>\\
def LoserTree<T>::Play(int p, int lc, int rc, int(*winner)(T A[], int b, int c),\\
int(*loser)(T A[], int b, int c)) \{\\
B[p] = loser(L, lc, rc);\\
int temp1, temp2;\\
（1）\underline{\hspace{12em}};\\
while p > 1 and p % 2: \{\\
temp2 = winner(L, temp1, B[p/2]);\\
（3）\underline{\hspace{12em}};\\
temp1 = temp2;\\
p/=2;\\
\}\\
（4）\underline{\hspace{12em}};\\
\}\\
3. （8 分）下述算法实现在图G 中计算从顶点i 到顶点j 之间长度为len 的简单路\\
径条数。图的ADT 如下：\\
Class Graph \{\\
public:\\
int VerticesNum();\\
int EdgesNum();\\
Edge FirstEdge(int oneVertex);\\
Edge NextEdge(Edge preEdge);\\
IsEdge(Edge onEdge);\\
int FromVertex(Edge oneEdge);\\
int ToVertex(Edge oneEdge);\\
\};\\
int visited[MAXSIZE];\\
// 初始化为0\\
int GetPathNum\_Len(G, int i, int j, int len)  \{\\
if (填空1 \underline{\hspace{12em}}) return 1;\\
sum = 0;\\
// sum 表示通过本结点的路径数\\
visited[i] = 1;\\
for (Edge e = G.FirstEdge(i); G.IsEdge(e); e = G.NextEdge(e)) \{\\
int v = G.ToVertex(e);\\
if (not visited[v])\\
填空2 \underline{\hspace{12em}}\\
\}// for\\
visited[i] = 0;\\
//本题允许曾经被访问过的结点出现在另一条路径中\\
return sum;\\
\} // GetPathNum\_Len\\
四、算法设计与分析 （25 分）\\
请尽量按照试题中的要求来写高效率的可读算法。应该申明算法思想，在代码\\
种加以恰当的注释。\\
1.\\
伸展树是一种自平衡的BST树，它可以通过旋转来实现树结构的动态调整。\\
伸展树结点的数据结构定义如下：\\
class TreeNode\{\\
int  data;\\
TreeNode  * father,* left,* right;\\
\};\\
其成员函数包括：\\
1）Delete zig(TreeNode* x); 其功能是删除以x 结点为根的子树；\\
2）Splay(TreeNode* x , TreeNode* f)；其功能是将 x 旋转为f 的子结点（f\\
是x 的祖先）。特殊的，把x 旋转到根结点即Splay(x, None)。\\
请运用上述给定的成员函数，实现删除树rt中所有大于u小于v的结点（假\\
设u, v是Splay树中的结点，且树种所有的节点值不重复）。\\
函数原型为：def DeleteSubTree(TreeNode* rt, TreeNode* u, TreeNode* v )\\
解答：把u 结点旋转到根；把v 旋转为u 的右儿子；删除v 结点的左子树前面几句\\
各3 分，最后一句1 分。如果思路、注释各1 分，不写则扣。\\
def DeleteUV(TreeNode* rt, TreeNode* u, TreeNode* v ) \{\\
Splay(u, None);\\
Splay(v, u);\\
Delete(v.lchild);\\
v.lchild = None;\\
\}\\
\\
\

\begin{solution}
本题为整卷转录型资料：选择题、填空题的参考答案以及简答、算法设计题的参考做法已统一放在本解答中，题面不保留原扫描夹带的答案标注。对扫描不清处，保留了原转录中可辨认的信息，后续可继续按原 PDF 图示补充图片。

\medskip
\noindent 原转录中夹带在题面里的参考答案/解答摘录如下，现已移入解答：
本卷包含选择填空、简答、算法设计等题型。下面保留按扫描件整理的题面；其中客观题参考答案已填入空格或括号，主观题给出参考思路或代码。

答案：\\
(1).根据题意，散列表长度为 L = 7/0.7 = 10；因此此题需要构建的哈希表是下标为\\
0\textasciitilde{}9 的一维数组。根据散列函数可以得到如下散列函数值表。\\
Key\\
H(Key)\\
采用线性探测法处理冲突，所构造的散列表为：\\
地址\\
关键字\\
（2）等概率情况下查找成功平均查找长度：\\
根据（1）的构造过程，查找次数如下表所示：\\
Key\\
Count\\
所以ASLsuccess= （1+1+1+1+3+3+2）/ 7 = 12/7。\\
等概率情况下查找不成功的平均查找长度：见散列表，计算查找不成功的次\\
数就直接找关键字到第一个地址上关键字为空的距离即可， 但根据哈希函数地址\\
为MOD7，因此初始只可能在0\textasciitilde{}6 的位置。等概率情况下，查找0\textasciitilde{}6 位置查找不\\
成功的次数表如下表所示\\
Key\\
Count\\
所以ASLunsuccess= （3+2+1+2+1+5+4）/ 7 = 18/7。\\
答案：\\
插入19 后结点持续分裂到根，整个树高增1，读盘3 次，写盘7 次，共读写10 次。\\
结果为：\\
删除48 时需先与后继50 交换，后删除，删除后各结点依然满足阶的要求，故不再\\
调整，删除过程中写盘2 次，若假设之前操作读进的页块不再读的话，读盘2 次，\\
或者从根开始重新读盘则4 次，结果如下：\\
答案：最终得到的红黑树：（用圆圈表示红色结点，方框表示黑色结点，外部空叶\\
结点省略不画）\\
答案：head(tail(head(tail(LS))))\\
答案：答案：\\
采用数学归纳法。\\
1. 当\\
\\
k\\
时，命题现任成立（只有一个点的情况）\\
2. 假设，当\\
n\\
k \\
时命题成立（即依次插入\\
1 \\
\\
n\\
个关键码递增的结点构成一棵\\
高度为n 的完全满二叉树），下面讨论\\
\\
n\\
k\\
的情况。\\
由于关键码的插入顺序单调递增，因此每次插入的新结点的初始位置一定是最\\
右链的叶子的右儿子。\\
按照插入顺序，当插入到关键码\\
1 \\
\\
n\\
时，构成一棵高度为n 的完全满二叉\\
O\\
E\\
S\\
A\\
I\\
Q\\
U\\
T\\
Y\\
N\\
树（归纳假设）。接着，考虑从关键码\\
2 \\
n\\
到\\
\\
\\
\\
\\
n\\
n\\
（共\\
2 \\
n\\
个关键码），\\
它们和根结点的右子树构成了\\
1 \\
\\
n\\
个连续递增的关键码，根据归纳假设，插入\\
关键码\\
\\
\\
\\
\\
n\\
n\\
后，根结点的右子树成为一棵高度为n 的完全满二叉树。然\\
后，下一个关键码\\
\\
\\
\\
\\
n\\
n\\
破坏了根结点的平衡（右子树比左子树深度大2），\\
根据AVL 树的旋转要求，树将以\\
n\\
2 为根（之后根结点不再参与旋转，因此这就是\\
最终的根结点），左子树是一棵高度为n 的完全满二叉树；而右子树几乎是一棵高\\
度为\\
\\
n\\
的完全满二叉树，唯一差别在于它的最右边还挂着\\
\\
\\
\\
\\
n\\
n\\
这个结\\
点。最后，关键码\\
\\
\\
\\
\\
n\\
n\\
到\\
2 \\
n\\
插入到树中，根据归纳假设，它们将把\\
右子树构造成一棵高度为n 的完全满二叉树。此时，整棵树便是一棵高度为\\
\\
n\\
的\\
完全满二叉树。\\
综上，结论成立。\\

\end{solution}
\end{question}
